% Standalone editable source. Compile with latexmk -pdf heide_math_v3.tex.
% No external figures, bibliography database, or shell escape required.
\documentclass[11pt,a4paper]{article}
\usepackage[T1]{fontenc}
\usepackage[utf8]{inputenc}
\usepackage{lmodern}
\usepackage[margin=25mm,headheight=14pt]{geometry}
\usepackage{amsmath,amssymb,bm}
\usepackage{booktabs,array,tabularx}
\usepackage{enumitem,microtype,xcolor,fancyhdr,needspace}
\usepackage[colorlinks=true,linkcolor=blue!45!black,citecolor=blue!45!black,
 urlcolor=blue!45!black,breaklinks=true]{hyperref}
\usepackage{bookmark}
\hypersetup{pdftitle={From Heide's Force Theorem to Magnetic Tensors},
 pdfauthor={DMI extraction project},pdfsubject={Mathematical foundation for spiralization, stiffness and exchange from spin spirals}}
\numberwithin{equation}{section}
\setcounter{tocdepth}{1}
\setlength{\parindent}{0pt}
\setlength{\parskip}{5pt plus 1pt minus 1pt}
\setlength{\emergencystretch}{2em}
\setlist{nosep,leftmargin=1.5em,topsep=5pt}
\renewcommand{\arraystretch}{1.17}
\newcolumntype{Y}{>{\raggedright\arraybackslash}X}
\pagestyle{fancy}
\fancyhf{}
\fancyhead[L]{\small Spin-spiral tensors from DFT}
\fancyhead[R]{\small Revision 3.0}
\fancyfoot[C]{\thepage}
\renewcommand{\headrulewidth}{0.3pt}
\newcommand{\vct}[1]{\bm{#1}}
\newcommand{\nh}{\hat{\vct n}}
\newcommand{\qh}{\vct q}
\newcommand{\kh}{\vct k}
\newcommand{\sh}{\hat{\vct s}}
\newcommand{\Nm}{N_{\mathrm m}}
\newcommand{\ii}{\mathrm i}
\newcommand{\Tr}{\operatorname{Tr}}
\newcommand{\Dt}{\mathcal D}
\newcommand{\Pj}{\mathcal P}
\newcommand{\code}[1]{\texttt{\detokenize{#1}}}

\begin{document}
\hypersetup{pageanchor=false}
\begin{titlepage}
\vspace*{15mm}
{\large Spin-spiral tensor pipeline\par}
\vspace{12mm}
{\LARGE\bfseries From Heide's Force Theorem\\[3pt] to Magnetic Tensors\par}
\vspace{8mm}
{\large Computing spiralization, stiffness and exchange\\
from paired, SOC-free spin spirals\par}
\vspace{10mm}
{\large Revision 3.0 \quad | \quad 2 November 2026\par}
\vspace{13mm}
\begin{abstract}
\noindent
The calculation has four parts. First, VASP prepares an unconstrained
SOC-free flat spin spiral for a chosen magnetic pattern. Second, a
frozen-state PAW trace gives its first-order spin--orbit energy, as a vector
with one component per Cartesian axis. Third, symmetry fixes which tensor
components are allowed and which $q$ points are needed. Fourth, small linear
fits turn traces and total energies into the spiralization tensor, the
stiffness tensor and exchange classes.

This document derives what the code needs for those steps and nothing more.
The tensors are long-wavelength sums of the bond interactions, so they avoid
the bond-resolution limit of flat spirals. They then give predictions such as
the spiral period, spin plane and texture class.
\end{abstract}
\vfill
\textbf{Source and scope.} This revision replaces the bond-vector target of
Revision 2.4. It keeps the generalized-Bloch states, Patch B frozen-state
trace and lab-frame moment readback. It removes the constrained families F1 and
F2, the bond basis and the paired bond rows. The new equations were checked
numerically on spin models and repository fixtures. No VASP calculation has
yet tested this revision.

\textbf{Build.} Run \code{latexmk -pdf heide_math_v3.tex}, or run
\code{pdflatex} twice. All text, equations and references are in this file.
\end{titlepage}
\hypersetup{pageanchor=true}
\pagenumbering{roman}
\tableofcontents
\clearpage
\pagenumbering{arabic}

\section{What the calculation measures}\label{sec:overview}

We consider a magnet with one slowly varying order parameter: a ferromagnet,
or a collinear pattern of site signs $\varepsilon_a=\pm1$ in the magnetic
primitive cell. A flat spiral twists that pattern with wavevector $\qh$ about
a spin-rotation axis $\nh$:
\begin{equation}\label{eq:flat}
 \hat{\vct S}_a(\vct r)=\varepsilon_a
 \bigl[\hat{\vct e}_1\cos(\qh\cdot\vct r)
 +\hat{\vct e}_2\sin(\qh\cdot\vct r)\bigr],
 \qquad \hat{\vct e}_1\times\hat{\vct e}_2=\nh.
\end{equation}
For long wavelengths, the energy per magnetic site relative to $\qh=0$ is
\begin{equation}\label{eq:energy}
 \boxed{e(\qh,\nh)=\qh^{T}A\,\qh+\nh^{T}\Dt\,\qh+O(q^3).}
\end{equation}
$A$ is the symmetric \emph{stiffness tensor}; it comes from SOC-free exchange.
$\Dt$ is the $3\times3$ \emph{spiralization tensor}; it comes from SOC at first
order. Anisotropy is second order in SOC and is treated separately in
Section~\ref{sec:predict}.

Two measured quantities isolate these tensors. Let $T_\alpha(\qh)$ be the
first-order SOC energy per calculation cell with the operator selecting
Cartesian axis $\alpha$, and let $E(\qh)$ be the SOC-free total energy per
cell. Then
\begin{align}
 y_\alpha(\qh)&=\frac{T_\alpha(\qh)-T_\alpha(-\qh)}{\Nm}
 =2\,\hat{\vct e}_\alpha^{\,T}\Dt\,\qh+O(q^3),\label{eq:y}\\
 g(\qh)&=\frac{E(\qh)+E(-\qh)-2E(0)}{2\Nm}
 =\qh^{T}A\,\qh+O(q^4).\label{eq:g}
\end{align}
$\Nm$ is the number of magnetic sites in the calculation cell. Note the factor
of two in Eq.~\eqref{eq:y}. It is the same paired observable as Revision 2.4,
with \emph{no additional factor of one-half}.

The same energies at larger $q$ also determine exchange classes $J_c$
(Section~\ref{sec:classes}). One least-squares routine handles all three
quantities; only the symmetry rule and the row differ.

Throughout, $\qh_{\rm frac}$ denotes reciprocal fractional coordinates and
$\vct r_{\rm frac}$ direct fractional coordinates:
$\qh\cdot\vct r=2\pi\qh_{\rm frac}\cdot\vct r_{\rm frac}$.
Without a subscript, $\qh$ is Cartesian in \AA$^{-1}$.
Energies are per magnetic site unless stated. $\Dt$ then has units of
eV\,\AA{} and $A$ of eV\,\AA$^2$.

\section{The DFT state and the first-order SOC energy}\label{sec:dft}

\subsection{What is frozen}

The electronic calculation follows the first-order spin-spiral method of
Heide, Bihlmayer and Bl\"ugel~\cite{heide}. For each spiral, VASP solves the
noncollinear SOC-free Kohn--Sham problem,
\begin{equation}\label{eq:ks}
 H^{(0)}_{\rm KS}[\rho,\vct m]\,\Psi^{(0)}_{n\kh}
 =\varepsilon^{(0)}_{n\kh}\Psi^{(0)}_{n\kh}.
\end{equation}
The first-order SOC insertion is
\begin{equation}\label{eq:firstorder}
 \Delta E^{(1)}=
 \sum_{n\kh}w_{\kh}f^{(0)}_{n\kh}
 \langle\Psi^{(0)}_{n\kh}|H_{\rm SOC}|\Psi^{(0)}_{n\kh}\rangle.
\end{equation}
This is the derivative of the stationary electronic energy with respect to a
small SOC coupling. Stationarity is why it can be evaluated on the unperturbed
orbitals and occupations. It is an SOC expectation, not VASP's full
\code{TOTEN}, so the two kinds of energy are never added together.

\subsection{Why a spin spiral fits in the chemical cell}

Without SOC, the electronic functional is invariant under a common spin
rotation, so a translation can be accompanied by a spin rotation. VASP's
\code{LNONCOLLINEAR = T}, \code{LSPIRAL = T} representation uses
\begin{equation}\label{eq:gb}
 \Psi_{n\kh}(\vct r)=
 \begin{pmatrix}
 e^{\ii(\kh-\qh/2)\cdot\vct r}\,u_{n\kh\uparrow}(\vct r)\\
 e^{\ii(\kh+\qh/2)\cdot\vct r}\,u_{n\kh\downarrow}(\vct r)
 \end{pmatrix},
 \qquad u(\vct r+\vct R)=u(\vct r).
\end{equation}
The transverse magnetization rotates by $\qh\cdot\vct R$ between cells. This
generalized Bloch construction avoids a supercell for every
\code{QSPIRAL}~\cite{spiralwiki}. SOC breaks this freedom, so the spiral is
prepared with \code{LSORBIT = F} and SOC is inserted afterward.

\subsection{Why the SCF stays unconstrained}\label{sec:unconstrained}

Each sublattice phase may differ from the rigid value $\qh\cdot\vct\tau_a$ by
a small amount $\delta_a$. Let $\xi$ scale the SOC. Then
\begin{equation}\label{eq:stationary}
 E(\xi,\vct\delta)=E_0(\vct\delta)+\xi\,T(\vct\delta)+O(\xi^2).
\end{equation}
Minimizing over $\vct\delta$ changes the energy only at $O(\xi^2)$, because
$E_0$ is stationary at its own minimum $\vct\delta^*$. The complete first-order
result is therefore $T(\vct\delta^*)$, evaluated on the \emph{relaxed} SOC-free
state. A rigid state generally has the wrong $E_0$ at $O(q^2)$ and, since
$\vct\delta^*\propto\qh$ when allowed, the wrong $T$ at $O(q)$.

We therefore run with \code{LAMBDA = 0} and let the SCF find the phases.
Symmetry decides when the relaxation can occur (Section~\ref{sec:twist}).
A flat spiral of one order parameter is already a stationary SOC-free state, so
the constrained-family machinery of Revision 2.4 is not needed.

\section{Preparing and reading one spiral}\label{sec:frames}

\subsection{The target}

Use \code{SAXIS = 0 0 1}. Lab frame means these fixed Cartesian spin axes,
independent of $\qh$. In the production states the rotation axis is always
$\hat{\vct z}$. For magnetic site $a$ at direct position $\vct\tau_{a}$,
\begin{equation}\label{eq:target}
 \boxed{\hat{\vct S}_a^0(\qh)=
 \varepsilon_a\bigl(\cos\phi_a,\ \sin\phi_a,\ 0\bigr),
 \qquad \phi_a=2\pi\qh_{\rm frac}\cdot\vct\tau_{a,\rm frac}.}
\end{equation}
This single formula covers every site count. The partner is the same formula
at $-\qh$, and the $q=0$ reference is the same formula at $\qh=0$.

\Needspace{70mm}
\subsection{VASP inputs}

\begin{center}
\begin{tabularx}{\textwidth}{@{}lY@{}}
\toprule Tag & Value and reason\\\midrule
\code{LNONCOLLINEAR}, \code{LSPIRAL} & \code{T}, also for the $q=0$ reference.\\
\code{LSORBIT} & \code{F}. SOC enters only through the trace.\\
\code{QSPIRAL} & $\qh_{\rm frac}$, in reciprocal lattice units.\\
\code{LZEROZ}, \code{SAXIS} & \code{T} and \code{0 0 1}: flat spiral about $z$.\\
\code{ISYM} & $-1$. VASP symmetry does not account for the spiral~\cite{spiralhowto}.\\
\code{I_CONSTRAINED_M}, \code{LAMBDA} & $1$ and $0$. This gives sphere moments with no penalty.\\
\code{RWIGS} & One explicit radius per species.\\
\code{ENCUT}, \code{ENINI} & \code{ENCUT} about 100 eV above \code{ENINI}, as the wiki advises.\\
\code{MAGMOM} & $m_a\hat{\vct S}_a^0$, with seed magnitude $m_a$.\\
\code{M_CONSTR} & $\hat{\vct S}_a^0$. It has no effect at \code{LAMBDA = 0}.\\
\bottomrule
\end{tabularx}
\end{center}
\textbf{Do not pre-unwind \code{MAGMOM}.} VASP's \code{COUNTER_SPIRAL} applies
the internal rotation $R_z(-\phi_a)$ itself~\cite{fetest}. Keep the $k$ mesh,
cutoffs, smearing and spheres identical across all $q$ of one material.

\subsection{Reading the state}

Read the final complete \code{M_int} block of OSZICAR, in POSCAR ion order.
The vectors are lab components, and they are the only source of site directions.
Let $M_a(\qh)$ be the length of that vector on site $a$. In-plane directions
are conveniently complex numbers:
\begin{equation}\label{eq:complex}
 s_a=\frac{M^{\rm int}_{a,x}+\ii M^{\rm int}_{a,y}}{M_a},
 \qquad u_a=\varepsilon_ae^{\ii\phi_a}.
\end{equation}
SOC-free energies do not depend on a common rotation about $z$, so remove the
best common angle before comparing with the target:
\begin{equation}\label{eq:delta}
 \theta=\arg\sum_a u_a^*s_a,\qquad
 \delta_a=\arg\bigl(u_a^*s_ae^{-\ii\theta}\bigr).
\end{equation}
$\delta_a$ is the realized phase relaxation of site $a$. It keeps the sign of
the moment: a reversed site gives $\delta_a\simeq180^\circ$. In the historical
$\lambda=0$ bcc Fe run, sites seeded along $+x$ reached a $45.34^\circ$
difference against the $45^\circ$ F0 value~\cite{fetest}. This shows that
unconstrained phases do move, and sets a first estimate of their convergence
error.

The total energy is the final OUTCAR \code{free  energy   TOTEN}, which prints
more digits than OSZICAR. With $\lambda=0$ it contains no penalty energy.

\section{From the PAW expression to the trace vector}\label{sec:paw}

\subsection{Why the all-electron on-site matrix appears}

Define the projector coefficients and all-electron on-site SOC matrix by
\begin{equation}\label{eq:pawobjects}
 c^\mu_{i,n\kh s}=
 \langle\tilde p_i^\mu|\tilde\Psi_{n\kh s}\rangle,\qquad
 C^{\mu,ss'}_{ij}=
 \langle\phi_i^\mu|H_{\rm SOC}^{ss'}|\phi_j^\mu\rangle.
\end{equation}
Here $\mu$ labels an atom and $i,j$ its partial-wave channels. The PAW
expectation of a local operator is a smooth term plus an
all-electron-minus-pseudo on-site correction. Within the one-center expansion
the smooth part inside a sphere cancels its pseudo subtraction. Neglecting SOC
outside the spheres leaves the all-electron on-site matrix, which is VASP's
approximation~\cite{lsorbit}. The validated \code{CSO} dump supplies $C$
directly, so no radial integral is coded.

\subsection{Split the operator into three orbital matrices}

The SOC operator is linear in the Pauli matrices. On each atom,
\begin{equation}\label{eq:split}
 C=\sum_{\alpha=x,y,z}\sigma_\alpha\otimes\Lambda_\alpha,
\end{equation}
where each $\Lambda_\alpha$ is an orbital matrix. For the atomic form,
$\Lambda_\alpha=\xi L_\alpha$. Reading off the four spin blocks gives
\begin{equation}\label{eq:blocks}
 \begin{pmatrix}
 C_{\uparrow\uparrow}&C_{\uparrow\downarrow}\\
 C_{\downarrow\uparrow}&C_{\downarrow\downarrow}
 \end{pmatrix}=
 \begin{pmatrix}
 \Lambda_z&\Lambda_x-\ii\Lambda_y\\
 \Lambda_x+\ii\Lambda_y&-\Lambda_z
 \end{pmatrix},
\end{equation}
and therefore
\begin{equation}\label{eq:lambda}
 \boxed{\Lambda_x=\tfrac12\bigl(C_{\uparrow\downarrow}+C_{\downarrow\uparrow}\bigr),
 \quad
 \Lambda_y=\tfrac{1}{2\ii}\bigl(C_{\downarrow\uparrow}-C_{\uparrow\downarrow}\bigr),
 \quad
 \Lambda_z=\tfrac12\bigl(C_{\uparrow\uparrow}-C_{\downarrow\downarrow}\bigr).}
\end{equation}
One \code{CSO} dump at \code{SAXIS = 0 0 1} thus supplies all three. Keep all
four blocks when reading it: the $x$ and $y$ information lives in the flip
blocks. The block order in the file is a stored convention, which T3 checks.

\subsection{Selecting the rotation axis}

The production state rotates about $z$. Let $U$ be a spin rotation with
$U\sigma_zU^\dagger=\nh\cdot\vct\sigma$. Because the SOC-free energy is
unchanged by $U$, the state $U|\Psi\rangle$ is the same spiral with rotation
axis $\nh$. Its SOC expectation is
\begin{equation}\label{eq:rotate}
 \langle U\Psi|C|U\Psi\rangle=\langle\Psi|U^\dagger CU|\Psi\rangle,
 \qquad
 U^\dagger CU=\sigma_z\otimes(\nh\cdot\vct\Lambda)+\text{(flip terms)}.
\end{equation}
From Eq.~\eqref{eq:gb}, a translated projection of spin $s=\pm1$ acquires
$e^{\ii(\kh-s\qh/2)\cdot\vct R}$. Diagonal blocks have no remaining phase, while
flip blocks carry $e^{\pm\ii\qh\cdot\vct R}$. Their lattice average vanishes:
\begin{equation}\label{eq:average}
 \frac{1}{N_R}\sum_{\vct R}e^{\pm\ii\qh\cdot\vct R}=0,
 \qquad \qh\notin\{\hbox{reciprocal lattice vectors}\}.
\end{equation}
The flip \emph{contribution} vanishes, not the stored flip blocks. Only
$\sigma_z\otimes(\nh\cdot\vct\Lambda)$ remains, which is linear in $\nh$.

\subsection{The trace vector}

One state therefore gives a whole vector, with no second SCF and no rotation of
the moments:
\begin{equation}\label{eq:trace}
 \boxed{T(\nh)=\nh\cdot\vct T,\qquad
 T_\alpha=\sum_{\mu n\kh}w_{\kh} f^{(0)}_{n\kh}
 \sum_{s=\pm1}s\,[\vct c^\mu_{n\kh s}]^\dagger\Lambda_\alpha^\mu
 \,\vct c^\mu_{n\kh s}.}
\end{equation}
Summing over occupied states first gives the density-matrix form,
\begin{equation}\label{eq:rho}
 \rho^{\mu s}_{ij}=
 \sum_{n\kh}w_{\kh}f^{(0)}_{n\kh}
 c^\mu_{i,n\kh s}c^{\mu *}_{j,n\kh s},\qquad
 \boxed{T_\alpha=\sum_\mu
 \Tr\bigl[\Lambda_\alpha^\mu
 \bigl(\rho^{\mu\uparrow}-\rho^{\mu\downarrow}\bigr)\bigr].}
\end{equation}
This needs only the on-site \code{COCC} blocks, not the full \code{CPROJ}
arrays, which is a large storage saving. The stored convention matters:
$\code{COCC}[i,j]=\rho[j,i]$. The implementation therefore contracts
\code{sum(Lambda * (COCC_uu - COCC_dd))} elementwise, not
\code{trace(Lambda @ ...)}, which inserts the wrong transpose. A wrong
transpose can still give a real number, so the imaginary-part gate cannot
replace this check.

The per-atom terms of Eq.~\eqref{eq:rho} are retained. They are additive,
so they later give the element-resolved origin of $\Dt$.

Equation~\eqref{eq:trace} equals the rotated-operator trace of Revision 2.4,
and the two agreed to $10^{-14}$ on repository fixtures for $x$, $y$, $z$ and
generic axes. It removes the Euler conventions and $U$ matrices from production.

\subsection{Parity and the null at \(q=0\)}\label{sec:parity}

Take a flat spiral with its common phase removed. A spin-only $180^\circ$
rotation about the in-plane $x$ axis sends $\phi\mapsto-\phi$, so it maps the
state at $\qh$ onto the state at $-\qh$. It also leaves the SOC-free energy
unchanged and changes $\sigma_z$ to $-\sigma_z$. Therefore, for every axis $\alpha$,
\begin{equation}\label{eq:parity}
 T_\alpha(-\qh)=-T_\alpha(\qh),\qquad E(-\qh)=E(\qh).
\end{equation}
Equivalently, $(\qh,\nh)$ and $(-\qh,-\nh)$ describe the same spiral.

At $\qh=0$ the state is collinear, and the same rotation about the collinear
axis maps it onto itself. Therefore $T_\alpha(0)=0$ for all three axes. This
also shows why $T$ contains no anisotropy.

These are exact statements about ideal states. Two independent SCFs may still
differ through their branch, occupations or convergence, so production computes
both signs and tests Eq.~\eqref{eq:parity} instead of assuming it.

\section{Which VASP run supplies each quantity}\label{sec:inputs}

The operator is a property of the material. The projections and occupations
belong to one state. This separation is the central Patch B input
rule~\cite{projectcore}.

\Needspace{55mm}
\begin{center}
\begin{tabularx}{\textwidth}{@{}lY@{}}
\toprule Run & Object and relevant condition\\\midrule
\code{dump_C} & Full \code{CSO} with SOC on and \code{LSPIRAL = F}, from a
converged SOC-free in-plane reference. One per material.\\
Genuine $x$, $y$ \code{dump_C} & \code{CSO} at \code{SAXIS} along $x$ and $y$.
Used only for check T3, once per build.\\
Spiral SCF & SOC-free state at $\pm\qh$ and at $\qh=0$. Supplies
$E$ and \code{M_int}.\\
\code{dump_P} & \code{COCC} from one SOC-free \code{PRE_LOOP} capture, after
\code{PROALL} and before orbital updates. Add \code{CPROJ} and \code{FERWE} on
the first state of each material for check M1.\\
Exchange SCF & SOC-free state on a coarse $q$ grid. No dump, no patch.\\
Anisotropy run & Standard SOC static run at $q=0$. No patch.\\
\bottomrule
\end{tabularx}
\end{center}
A dump stage is not a substitute for a converged parent SCF. A common
\code{CSO} is an approximation across states, because its radial factors depend
on the potential. This is expected to be small for small $q$, and the
supercell anchor in Section~\ref{sec:checks} tests it.

\section{What the tensors mean in a spin model}\label{sec:model}

The DFT quantities do not need a spin model. The model below gives the tensors
their meaning, and it is the oracle for the code tests.

Use the Hamiltonian of Revision 2.4, with each undirected bond $b$ counted once:
\begin{equation}\label{eq:ham}
 H=\sum_b\bigl[-J_b\,\hat{\vct S}_i\cdot\hat{\vct S}_j
 +\vct d_b\cdot(\hat{\vct S}_i\times\hat{\vct S}_j)\bigr].
\end{equation}
A bond runs from site $a$ in the home cell to site $a'$ at displacement
$\vct r_b$. Write $\varepsilon_b=\varepsilon_a\varepsilon_{a'}$. On the flat
spiral of Eq.~\eqref{eq:flat}, the two ends differ in phase by
$\qh\cdot\vct r_b$, so
\begin{equation}\label{eq:modelgy}
 g(\qh)=\frac{1}{\Nm}\sum_b\varepsilon_bJ_b\bigl[1-\cos(\qh\cdot\vct r_b)\bigr],
 \qquad
 y_\alpha(\qh)=\frac{2}{\Nm}\sum_b\varepsilon_b\,d_{b,\alpha}\sin(\qh\cdot\vct r_b).
\end{equation}
Expanding at small $q$ and comparing with Eqs.~\eqref{eq:y} and~\eqref{eq:g} gives
\begin{equation}\label{eq:tensors}
 \boxed{\Dt=\frac{1}{\Nm}\sum_b\varepsilon_b\,\vct d_b\,\vct r_b^{\,T},
 \qquad
 A=\frac{1}{2\Nm}\sum_b\varepsilon_bJ_b\,\vct r_b\,\vct r_b^{\,T}.}
\end{equation}
Reversing a bond reverses both $\vct r_b$ and $\vct d_b$, so each term is
independent of the chosen orientation. Equation~\eqref{eq:tensors} shows why
this pivot works. A flat spiral cannot separate two bonds with the same
displacement class, but it does not need to: the tensors are exactly the sums it
can measure.

\textbf{A normalization check.} Take a chain with one site per cell, spacing
$a$, one bond per site, exchange $J$ and $\vct d=d\,\hat{\vct z}$. The only
nonzero tensor entries for this chain are $\Dt_{zx}=da$ and $A_{xx}=Ja^2/2$.
Equation~\eqref{eq:modelgy} gives $y_z=2d\sin(qa)\simeq2\Dt_{zx}q$ and
$g=J(1-\cos qa)\simeq A_{xx}q^2$, as required.

\section{Symmetry}\label{sec:symmetry}

Symmetry is used three ways: to count parameters before any VASP run, to choose
informative $q$ directions, and to test the result. Only the crystal structure
and the signs $\varepsilon_a$ are needed.

\subsection{Operations in Cartesian form}

\code{spglib}~\cite{spglib} returns operations $(W_g,\vct t_g)$ acting on direct
fractional coordinates, $\vct x\mapsto W_g\vct x+\vct t_g$. Let the columns of $L$ be the
lattice vectors, so that $L$ is the transpose of the POSCAR matrix. Then
\begin{equation}\label{eq:cart}
 \boxed{R_g=LW_gL^{-1}.}
\end{equation}
The tensors use only $R_g$. The translation $\vct t_g$ matters for the site
permutation $\pi_g$, defined by
\begin{equation}\label{eq:perm}
 W_g\vct\tau_{a}+\vct t_g\equiv\vct\tau_{\pi_g(a)}
 \quad(\text{modulo lattice vectors}),
\end{equation}
with positions in fractional coordinates. It maps magnetic sites to magnetic sites.

\subsection{Which operations apply}

An operation is kept if it maps the sign pattern to itself or to its overall
reverse. That is, there must be one sign $\eta_g=\pm1$ with
\begin{equation}\label{eq:keep}
 \boxed{\varepsilon_{\pi_g(a)}=\eta_g\,\varepsilon_a\quad\text{for every }a.}
\end{equation}
Call the kept set $G$. If this test fails, the operation connects two different
magnetic patterns and says nothing about the coefficients of this one.
Magnetic space groups are not needed. They describe one fixed spin orientation,
while $\Dt$ and $A$ describe the energy for all orientations of the pattern.

\subsection{The transformation rule}

Apply a kept operation to the crystal and the spiral together. The crystal is
unchanged, positions and $\qh$ are polar, and spin is axial:
\begin{equation}\label{eq:maps}
 \qh\mapsto R\qh,\qquad
 \hat{\vct S}\mapsto\det(R)\,R\,\hat{\vct S},\qquad
 \nh\mapsto\det(R)\,R\,\nh.
\end{equation}
The last follows because $\nh$ is a cross product of two spins. If $\eta_g=-1$,
every spin also reverses. That shifts the spiral phase by $\pi$ and leaves $\nh$
and the energy unchanged, which is why sublattice-exchanging operations still
constrain the tensors. Hence $e(R\qh,\det(R)R\nh)=e(\qh,\nh)$. Since each
operation's inverse is also kept, Eq.~\eqref{eq:energy} gives
\begin{equation}\label{eq:rules}
 \boxed{\Dt=\det(R)\,R\,\Dt\,R^{T},\qquad A=A^T=R\,A\,R^{T}\qquad\text{for all }g\in G.}
\end{equation}
Without the $\det(R)$ the rule is wrong for every improper operation. If
inversion is kept, then $\Dt=-\Dt$, so $\Dt=0$.

\subsection{Basis by group average}\label{sec:basis}

Averaging over the group projects any matrix onto the allowed form:
\begin{equation}\label{eq:project}
 \boxed{\Pj_{\Dt}[M]=\frac{1}{|G|}\sum_{g\in G}\det(R_g)\,R_gMR_g^{T},
 \qquad
 \Pj_{A}[M]=\frac{1}{|G|}\sum_{g\in G}R_gMR_g^{T}.}
\end{equation}
Apply $\Pj_{\Dt}$ to the nine unit matrices, and $\Pj_{A}$ to the six symmetric
unit matrices. A singular-value decomposition of the results gives orthonormal
bases, so that
\begin{equation}\label{eq:basis}
 \Dt=\sum_{\ell=1}^{p_{\Dt}}d_\ell B_\ell,\qquad
 A=\sum_{\ell=1}^{p_A}a_\ell C_\ell.
\end{equation}
The rank of each set is the number of independent parameters, $p_{\Dt}$ or
$p_A$. The same projector also gives the forbidden part of a measured tensor
in Section~\ref{sec:fit}.

These results are useful regression tests. In the table, $\hat{\vct z}$ is the
principal axis, $[\hat{\vct z}]_\times$ is the matrix of
$\vct v\mapsto\hat{\vct z}\times\vct v$, and in-plane stiffness is
$A_\perp$ with axial stiffness $A_\parallel$.

\Needspace{50mm}
\begin{center}
\begin{tabularx}{\textwidth}{@{}lllcY@{}}
\toprule Point group & $\Dt$ & $A$ & $p_{\Dt},p_A$ & Texture class\\\midrule
Inversion kept & $0$ & depends on the group & $0,\ge1$ & None\\
$23$ & $D\,I$ & $A\,I$ & $1,1$ & Bloch\\
$\bar43m$ & $0$ & $A\,I$ & $0,1$ & None\\
$3m$, $4mm$, $6mm$ & $D[\hat{\vct z}]_\times$ & $\mathrm{diag}(A_\perp,A_\perp,A_\parallel)$ & $1,2$ & N\'eel\\
$622$ & $\mathrm{diag}(D_\perp,D_\perp,D_\parallel)$ & $\mathrm{diag}(A_\perp,A_\perp,A_\parallel)$ & $2,2$ & Bloch\\
$\bar42m$, 2-fold axes on $x,y$ & $D\,\mathrm{diag}(1,-1,0)$ & $\mathrm{diag}(A_\perp,A_\perp,A_\parallel)$ & $1,2$ & Antiskyrmion\\
\bottomrule
\end{tabularx}
\end{center}

\subsection{Intracell twist: when rigid spirals fail}\label{sec:twist}

Write the relaxed phases as $\phi_a=\qh\cdot\vct\tau_a+\delta_a(\qh)$. For small
$q$, $\delta_a=\vct\eta_a\cdot\qh$, with one vector $\vct\eta_a$ per magnetic
site. A common shift of all $\vct\eta_a$ only moves the spiral's origin, so set
$\sum_a\vct\eta_a=0$.

In the SOC-free state a spatial operation moves sites but not spins. It maps the
relaxed spiral at $\qh$ to the relaxed spiral at $R\qh$, with site $a$ becoming
site $\pi(a)$. This requires $\vct\eta_{\pi(a)}=R\vct\eta_a+\vct s_g$ for some
common $\vct s_g$. Summing over $a$ gives $\vct s_g=0$, so
\begin{equation}\label{eq:twist}
 \vct\eta_{\pi_g(a)}=R_g\vct\eta_a\quad\text{for all }g\in G,\qquad
 \sum_a\vct\eta_a=0.
\end{equation}
Collect the $\vct\eta_a$ into one vector of length $3\Nm$. Let $\mathsf G_g$ be
the matrix that places block $R_g$ at row $\pi_g(a)$ and column $a$, and let
$\mathsf Q$ remove the common shift. The number of allowed twist parameters is
\begin{equation}\label{eq:ntwist}
 \boxed{n_{\rm twist}=\operatorname{rank}\Bigl[
 \frac{1}{|G|}\sum_{g\in G}\mathsf G_g\,\mathsf Q\Bigr],
 \qquad
 \mathsf Q=I-\frac{1}{\Nm}\vct 1\vct 1^{T}\otimes I_3.}
\end{equation}

If $n_{\rm twist}=0$, rigid and relaxed spirals agree at first order in $q$. If
$n_{\rm twist}>0$, the rigid spiral is biased. An FeGe-like toy model with
$n_{\rm twist}=1$ shifted $A$ by $7\%$ and $\Dt$ by $35\%$; a BFO-like model with
$n_{\rm twist}=0$ showed no change with independent $J$ on every bond orbit. The
FeGe numbers come from arbitrary model couplings and show only the possible size
of the effect. A count of one for B20 agrees with the single twist function of
the published FeGe analysis~\cite{twisted}.

\textbf{A twist check.} A dimerized antiferromagnetic chain with period $a$,
bond lengths $x$ and $a-x$, and couplings $J_1$ and $J_2$ has
\begin{equation}\label{eq:chain}
 A_{\rm rigid}=\frac{|J_1|x^2+|J_2|(a-x)^2}{4},
 \qquad
 A_{\rm relaxed}=\frac{a^2}{4}\,\frac{|J_1J_2|}{|J_1|+|J_2|}.
\end{equation}
The relaxed form is the series-spring result. Rigid and relaxed agree only when
the two bond phases are already balanced, at $|J_1|x=|J_2|(a-x)$.

\subsection{Exchange classes}\label{sec:classes}

For exchange, the data determine one coupling per class of equivalent
displacements. List the bonds from all site pairs and cell translations within
a cutoff $R_{\max}$, counting each undirected bond once. Two bonds belong to the same class $c$ if they have equal
$\varepsilon_b$ and $\vct r'=\pm R_g\vct r$ for some $g\in G$. Equation~\eqref{eq:modelgy}
then becomes linear in the class couplings:
\begin{equation}\label{eq:classes}
 \boxed{g(\qh)=\sum_cJ_c\,\Phi_c(\qh),\qquad
 \Phi_c(\qh)=\frac{1}{\Nm}\sum_{b\in c}\varepsilon_b
 \bigl[1-\cos(\qh\cdot\vct r_b)\bigr].}
\end{equation}
Bonds that share a displacement class merge into one effective $J_c$, in
keeping with the flat-spiral limit. Use the true Cartesian $\qh$ in the row, not
$\qh$ reduced to the first zone: for sites inside the cell, $\qh$ and
$\qh+\vct G$ are different F0 states.

\section{Choosing the \(q\) points}\label{sec:qplan}

\subsection{Conversions}

With $L$ as above,
\begin{equation}\label{eq:qconvert}
 \qh=2\pi L^{-T}\qh_{\rm frac},\qquad
 \qh_{\rm frac}=\frac{1}{2\pi}L^{T}\qh.
\end{equation}
A measurement point is a unit direction $\sh$ and a magnitude $t$, with
$\qh=t\sh$. Each one requires the pair $\pm\qh$.

\subsection{Directions}

Two rank conditions must both hold. They are different, and the code checks both
before any input is written.

\begin{itemize}
\item \textbf{For $\Dt$:} let $S$ have the directions $\sh_i$ as columns. Each
direction gives three numbers, so three non-coplanar directions determine all
nine entries with no symmetry input. This requires $\operatorname{rank}S=3$.
\item \textbf{For $A$:} each direction gives only one number,
$\sh^TA\sh$. The matrix with entries $\sh_i^TC_\ell\sh_i$ must have rank $p_A$.
Six directions are needed with no symmetry; fewer suffice when $p_A$ is small.
\end{itemize}

Symmetry-equivalent directions can pass the first test and fail the second. In
$R3c$ BFO, the three directions $(100)$, $(010)$, $(001)$ in $\qh_{\rm frac}$ are
related by the threefold axis, so they measure the same stiffness. They give rank
1 of 2. The set $(111)$, $(1\bar10)$, $(100)$ gives rank 2 and $\operatorname{rank}S=3$.
For a polar crystal a good default is along the axis, in the basal plane, and
one oblique direction. The direction along the axis also gives $y=0$ by symmetry,
which is a direct test.

Choose from low-index candidates, and take the set with the largest smallest
singular value in both conditions. Use at least $\max(3,p_A)$ directions.

\subsection{Magnitudes}

Use three magnitudes per direction, so that one degree of freedom remains after
fitting two terms. The stiffness signal grows as $t^2$ and the spiralization
signal as $t$, so the stiffness sets the lower bound on $t$. The series truncation
sets the upper bound, roughly $t\,r_{\rm max}\lesssim1$ for the important
interaction range. For BFO-scale values, a first choice is
$t=0.03,\,0.06,\,0.09$ \AA$^{-1}$, to be confirmed by convergence tests.

\subsection{The \(k\) mesh and the exchange grid}

Use one $\Gamma$-centered mesh for every $q$ of a material, including $q=0$. Insulators
do not require $q$ to be commensurate with this mesh. Metals do need the mesh to
resolve $q$ near the Fermi surface, so use a convergence series there.

For exchange, take a coarse $\Gamma$-centered grid in $\qh_{\rm frac}$, remove
$q=0$, and keep one point per set related by $G$ and by $\qh\mapsto-\qh$.

\section{Fitting}\label{sec:fit}

Every fit below is ordinary linear least squares. The same small routine does
all of them.

\subsection{Slopes and curvatures along each direction}

For each direction $\sh$, fit the three magnitudes separately for each axis:
\begin{equation}\label{eq:slopes}
 y_\alpha(t\sh)=v_\alpha t+c_\alpha t^3,\qquad
 g(t\sh)=\kappa\,t^2+b\,t^4.
\end{equation}
The higher terms are nuisance parameters. The useful results are $\vct v(\sh)$
and $\kappa(\sh)$, which by Eqs.~\eqref{eq:y} and~\eqref{eq:g} are
\begin{equation}\label{eq:vkappa}
 \vct v(\sh)=2\Dt\,\sh,\qquad \kappa(\sh)=\sh^{T}A\,\sh.
\end{equation}

\subsection{The tensors}

Let $V$ have the vectors $\vct v(\sh_i)$ as columns, in the order of $S$. The
free spiralization tensor needs no symmetry input:
\begin{equation}\label{eq:free}
 \boxed{\Dt_{\rm free}=\tfrac12\,V\,S^{+},}
\end{equation}
where $S^{+}$ is the pseudoinverse, equal to $S^{-1}$ for three directions.
Its forbidden part is
\begin{equation}\label{eq:forbidden}
 \mathcal F=\Dt_{\rm free}-\Pj_{\Dt}[\Dt_{\rm free}].
\end{equation}
This should vanish within error. An axis, frame or transpose error generally
makes it nonzero. An overall sign error does not, so the sign must come from an
external anchor.

The reported tensor uses the symmetry basis. Minimize
\begin{equation}\label{eq:dfit}
 \sum_{i,\alpha}\Bigl[v_\alpha(\sh_i)-\sum_\ell d_\ell\,
 2\,\hat{\vct e}_\alpha^{\,T}B_\ell\sh_i\Bigr]^2.
\end{equation}
For $A$, the basis is required unless six or more directions are available. Minimize
\begin{equation}\label{eq:afit}
 \sum_i\Bigl[\kappa(\sh_i)-\sum_\ell a_\ell\,\sh_i^{T}C_\ell\sh_i\Bigr]^2.
\end{equation}
The eigenvalues of $A$ must be positive. A negative eigenvalue means the chosen
pattern is not even a local minimum of the exchange energy.

\subsection{Exchange classes}

Stack Eq.~\eqref{eq:classes} over the exchange grid and the small-$q$ points:
\begin{equation}\label{eq:jfit}
 \min_{\{J_c\}}\sum_m\Bigl[g(\qh_m)-\sum_cJ_c\Phi_c(\qh_m)\Bigr]^2.
\end{equation}
Check the rank first. Raise the bond cutoff until the leave-one-out prediction
error stops improving. The result gives a second estimate of the stiffness,
\begin{equation}\label{eq:acheck}
 A=\frac{1}{2\Nm}\sum_cJ_c\sum_{b\in c}\varepsilon_b\,\vct r_b\vct r_b^{\,T},
\end{equation}
which must agree with Eq.~\eqref{eq:afit}.

\subsection{Uncertainties}

For one linear fit with design matrix $X$, $m$ rows and $p$ parameters,
\begin{equation}\label{eq:cov}
 \hat\sigma^2=\frac{\|\vct y-X\hat{\vct x}\|^2}{m-p},\qquad
 \operatorname{Cov}(\hat{\vct x})=\hat\sigma^2(X^TX)^{-1}.
\end{equation}
With three magnitudes and two terms each direction has one degree of freedom,
so pool $\hat\sigma^2$ over all directions and axes in Eq.~\eqref{eq:slopes}.
Carry the resulting slope errors into Eqs.~\eqref{eq:dfit} and~\eqref{eq:afit}.
Axes from one state share noise, so treat these errors as estimates. Also report
the parity residuals of Section~\ref{sec:parity}, since a common
background cancels in $y$ and would not appear in this fit.

\subsection{Element-resolved spiralization}

Equation~\eqref{eq:rho} is a sum over atoms, and Eq.~\eqref{eq:y} is linear in
$T$. Applying the same fit to the atom-resolved $T^\mu_\alpha$ gives
$\Dt=\sum_\mu\Dt^\mu$, which can be summed by species. It shows
where the SOC energy is collected, within the on-site approximation.

\section{Predicted properties}\label{sec:predict}

\subsection{The spiral}

Minimize Eq.~\eqref{eq:energy} over $\qh$, then over the unit vector $\nh$:
\begin{equation}\label{eq:qstar}
 \boxed{\begin{gathered}
 \Gamma=\Dt A^{-1}\Dt^{T},\qquad
 \nh^*=\text{top eigenvector of }\Gamma\ \text{(eigenvalue $\gamma_{\max}$)},\\
 \qh^*=-\tfrac12A^{-1}\Dt^{T}\nh^*,\qquad
 e^*=-\tfrac14\gamma_{\max},\qquad
 \lambda=\frac{2\pi}{|\qh^*|}.
 \end{gathered}}
\end{equation}
Three cases are worth keeping as tests.

\begin{itemize}
\item \textbf{Scalar case,} $\Dt=D\,I$, $A$ scalar: $\lambda=4\pi A/|D|$ and
$e^*=-D^2/4A$. Then $\qh^*\parallel\nh^*$, a Bloch helix.
\item \textbf{Polar case,} $\Dt=D[\hat{\vct p}]_\times$:
$\nh^*\perp\hat{\vct p}$, $\qh^*=(D/2A_\perp)\,\hat{\vct p}\times\nh^*$ and
$\lambda=4\pi A_\perp/|D|$. This is a cycloid whose spin plane contains
$\hat{\vct p}$. Every direction in the basal plane is degenerate at this order,
so $\Dt$ does not select the propagation direction.
\item \textbf{Chain,} from Section~\ref{sec:model}: $q^*=-d/(Ja)$ and
$e^*=-d^2/2J$. For $J=10$ meV, $d=1$ meV, $a=3$ \AA{} this gives
$\lambda=18.85$ nm, against $18.91$ nm from the exact lattice minimum.
\end{itemize}

The texture class follows from the form of $\Dt$, as in the table of
Section~\ref{sec:basis}. These are harmonic, homogeneous results. A real spiral
ground state also depends on anisotropy and magnetostatics.

\subsection{Exchange and magnons}

For a ferromagnet with moment $M$ in $\mu_{\rm B}$ per site, spin $S=M/2$
and a cone of small angle $\theta$ about spin direction $\hat{\vct m}$, the
spiral energy per site is $\theta^2e$ and the spin deviation per site
is $S\theta^2/2$. Their ratio is the magnon energy,
\begin{equation}\label{eq:magnon}
 \hbar\omega(\kh)=\frac{4}{M}\bigl[g(\kh)+\hat{\vct m}^{T}\Dt\,\kh\bigr],
 \qquad D_{\rm sw}=\frac{4A}{M}.
\end{equation}
Here $\kh$ is the propagation direction of the wave
$S^x+\ii S^y\propto e^{\ii(\kh\cdot\vct r-\omega t)}$, and $\hat{\vct m}$ follows
$\langle\vct\sigma\rangle$, as VASP's magnetization does. The physical moment is
opposite, which reverses the sign of the linear term when written in terms of it.
A direct spin-dynamics integration on a chain confirmed this sign and prefactor.

The mean-field ordering temperature uses the class constant
\begin{equation}\label{eq:tmf}
 k_{\rm B}T_{\rm MF}=\frac23\,C,\qquad
 C=\frac{1}{\Nm}\sum_b\varepsilon_bJ_b=\sum_cJ_c\Phi_c^{\infty},\qquad
 \Phi_c^{\infty}=\frac{1}{\Nm}\sum_{b\in c}\varepsilon_b.
\end{equation}
This applies when all magnetic sites are equivalent. Mean field overestimates
ordering temperatures.

\subsection{Field and anisotropy}

For a cubic helimagnet, a cone of angle $\theta$ gives
$\sin^2\!\theta\,e^*-M\mu_{\rm B}B\cos\theta$. The cone closes at the saturation field
\begin{equation}\label{eq:bd}
 B_D=\frac{D^2}{2AM\mu_{\rm B}}.
\end{equation}
Anisotropy is second order in SOC, so the trace cannot supply it. Take $K$ per
magnetic site as the difference between hard and easy directions, from standard SOC
total energies at $q=0$. For a chosen pair $(\sh,\nh)$, define $A_s=\sh^TA\sh$ and
$D_s=\nh^T\Dt\sh$. The one-dimensional energy per site is
\begin{equation}\label{eq:sg}
 e[\varphi]=A_s(\varphi')^2+D_s\varphi'+K\sin^2\varphi.
\end{equation}
For a uniaxial easy-axis magnet, a $180^\circ$ wall then costs
\begin{equation}\label{eq:wall}
 \sigma=4\sqrt{A_sK}-\pi|D_s|,
\end{equation}
which becomes an energy per unit wall area on dividing by $V_{\rm m}$ of
Section~\ref{sec:units}. The collinear state gives way to a spiral when $\sigma$ becomes negative:
\begin{equation}\label{eq:dc}
 |D_s|>D_c=\frac{4}{\pi}\sqrt{A_sK}.
\end{equation}

\subsection{Units}\label{sec:units}

Let $V_{\rm m}=V_{\rm cell}/\Nm$ be the volume per magnetic site. Dividing a
per-site quantity by $V_{\rm m}$ gives the usual micromagnetic parameter:
\begin{center}
\begin{tabular}{@{}lll@{}}
\toprule Per site & Divide by $V_{\rm m}$ in \AA$^3$ & Multiply by\\\midrule
$A$ in meV\,\AA$^2$ & meV\,\AA$^{-1}$ & $1.602$ for pJ\,m$^{-1}$\\
$D$ in meV\,\AA & meV\,\AA$^{-2}$ & $16.02$ for mJ\,m$^{-2}$\\
$K$ in meV & meV\,\AA$^{-3}$ & $160.2$ for MJ\,m$^{-3}$\\
\bottomrule
\end{tabular}
\end{center}
Published normalizations differ, for example per atom or per formula unit, so
compare ratios such as $\lambda$ first.

\section{Checks that make the result usable}\label{sec:checks}

Each check addresses a specific way the result can fail.

\begin{description}[style=nextline,leftmargin=0pt,labelsep=0pt]
\item[C0, T3 and M1: is the operator and trace assembled correctly?]
C0 checks $C_{\downarrow\downarrow}=-C_{\uparrow\uparrow}$, an imaginary
antisymmetric $C_{\uparrow\uparrow}$ with zero diagonal, and flip-block
Hermiticity. T3 compares $\Lambda_x$ and $\Lambda_y$ from Eq.~\eqref{eq:lambda}
with the $C_{\uparrow\uparrow}$ block of genuine dumps at \code{SAXIS} along $x$
and $y$. M1 compares \code{COCC} with the density matrix rebuilt from
\code{CPROJ}, to $10^{-10}$ eV.

\item[SCF and state: is this the intended spiral?]
Require the final electronic loop to reach \code{EDIFF}. Require small $M_z$, and
$|\delta_a|<1^\circ$ when $n_{\rm twist}=0$. Otherwise record $\delta_a$; they should be
linear in $q$. Require $M_a(\qh)$ within $5\%$ of $M_a(0)$.

\item[N0, I0, P0 and E0: does the trace have its exact properties?]
N0 requires $T_\alpha(0)$ below $\max(5\times10^{-5}\ {\rm eV},3\sigma)$ on all
three axes. I0 bounds the imaginary residual below $10^{-6}$ eV. P0 compares
$T_\alpha(\qh)+T_\alpha(-\qh)$ with $T_\alpha(\qh)-T_\alpha(-\qh)$; start at $10\%$
at the smallest $q$ and tighten from data. E0 compares $E(\qh)$ with $E(-\qh)$.

\item[Forbidden: do the tensors obey symmetry?]
Require $\mathcal F$ of Eq.~\eqref{eq:forbidden} to vanish within $3\sigma$. Also
test directions with a symmetry-required zero, such as $y$ along a polar axis.

\item[Cross-check: do the two stiffness routes agree?]
Compare Eq.~\eqref{eq:afit} with Eq.~\eqref{eq:acheck}.

\item[Nulls and anchors: is the result physically right?]
If inversion is kept, then $e(-\qh,\nh)=e(\qh,\nh)$, so $y_\alpha(\qh)=0$ at every
$q$, not only in the slope. Bcc Fe is the first such test. For sign and magnitude, compare with direct SOC energies on a
commensurate supercell of a cheap noncentrosymmetric control. For the method,
compare with FLEUR first-order spiral results for BFO~\cite{xu} and
FeGe~\cite{grytsiuk}, matching structure, functional and normalization. These are
tests of the code, not of experiment: those calculations give about 80~nm against
62~nm for BFO and about 140~nm against about 70~nm for FeGe.
\end{description}

The remaining physical assumptions are compact: one order parameter with one
magnetic pattern; SOC at first order on fixed SOC-free states and occupations;
all-electron on-site SOC; and one material operator reused across small $q$.
SOC-driven twist is second order in the energy and is outside this
scheme~\cite{twisted}. The numerical checks of this document confirm the algebra.
They do not establish any DFT result.

\section{One material, from inputs to tensors}\label{sec:walkthrough}

\begin{enumerate}[leftmargin=1.8em]
\item Read the structure and $\varepsilon_a$. Build $G$ from
Eqs.~\eqref{eq:cart}--\eqref{eq:keep}.
\item Compute $B_\ell$, $C_\ell$ and $n_{\rm twist}$ from
Eqs.~\eqref{eq:project}--\eqref{eq:ntwist}.
\item Choose directions and magnitudes, and test both rank conditions of
Section~\ref{sec:qplan}.
\item Write targets from Eq.~\eqref{eq:target} at $\pm\qh$, at $q=0$ and on the
exchange grid. Run the SCFs, the \code{dump_P} captures and the \code{dump_C} runs.
\item Read $E$ and \code{M_int}. Compute $\delta_a$ from Eq.~\eqref{eq:delta}.
\item Form $\Lambda_\alpha$ from Eq.~\eqref{eq:lambda} and $\vct T$ from
Eq.~\eqref{eq:rho} for each state.
\item Form $y_\alpha$ and $g$ from Eqs.~\eqref{eq:y} and~\eqref{eq:g}.
\item Fit Eq.~\eqref{eq:slopes}, then compute $\Dt_{\rm free}$, $\mathcal F$, $\Dt$,
$A$ and $J_c$.
\item Apply the checks of Section~\ref{sec:checks}.
\item Evaluate Section~\ref{sec:predict} and convert units. Report the tensors in
the POSCAR Cartesian frame, with the stated Hamiltonian and per-site convention.
\end{enumerate}

\renewcommand{\refname}{References and implementation records}
\begin{thebibliography}{10}
\raggedright
\addcontentsline{toc}{section}{References and implementation records}
\bibitem{heide}
M. Heide, G. Bihlmayer and S. Bl\"ugel,
\emph{Describing Dzyaloshinskii--Moriya spirals from first principles},
Physica B \textbf{404}, 2678--2683 (2009).
\href{https://juser.fz-juelich.de/record/5598/files/HBB_FCM08.pdf}{Author preprint}.

\bibitem{xu}
B. Xu, S. Meyer, M. J. Verstraete, L. Bellaiche and B. Dup\'e,
\emph{First-principles study of spin spirals in the multiferroic BiFeO$_3$},
Phys. Rev. B \textbf{103}, 214423 (2021).
\href{https://doi.org/10.1103/PhysRevB.103.214423}{DOI}.

\bibitem{grytsiuk}
S. Grytsiuk \emph{et al.},
\emph{Ab initio analysis of magnetic properties of the prototype B20 chiral magnet FeGe}.
\href{https://arxiv.org/abs/1909.13545}{arXiv:1909.13545}.

\bibitem{twisted}
\emph{Micromagnetic description of twisted spin spirals in the B20 chiral magnet
FeGe from first principles}.
\href{https://arxiv.org/abs/2103.09800}{arXiv:2103.09800}.

\bibitem{spiralwiki}
VASP Wiki, \href{https://vasp.at/wiki/Spin_spirals}{Spin spirals}.
Generalized Bloch functions and magnetization.

\bibitem{spiralhowto}
VASP Wiki, \href{https://vasp.at/wiki/Spin-spiral_calculations}{Spin-spiral calculations}.
Symmetry, cutoffs and sphere-moment readback.

\bibitem{lsorbit}
VASP Wiki, \href{https://vasp.at/wiki/LSORBIT}{LSORBIT}.
All-electron one-center SOC approximation.

\bibitem{spglib}
\href{https://spglib.readthedocs.io/en/v2.8.0/api/autodoc/spglib.spg.html}{spglib documentation}.
Symmetry operations in direct fractional coordinates.

\bibitem{fetest}
Project bcc Fe convention tests and user-supplied VASP source excerpts,
reviewed 20--25 September 2026, as recorded in Revision 2.4.
Executable hashes and full run archives were not supplied with these excerpts.

\bibitem{projectcore}
Revision 2.4, \emph{From Heide's Force Theorem to the VASP Trace},
26 September 2026, for the Patch B run separation, stored \code{COCC} convention,
gate definitions and lab-frame moment contract, which this revision keeps.
\end{thebibliography}
\end{document}
